\documentclass[10pt,a4paper]{amsart}
\usepackage[margin=19mm]{geometry}
\usepackage{amsmath,amssymb,mathtools}
\usepackage[scr=rsfs,cal=euler]{mathalfa}
\usepackage{xfrac}
\usepackage[unicode,hidelinks,bookmarks=false]{hyperref}
\newcommand{\ie}{i.e.\ }
\newcommand{\cf}{cf.\ }
\newcommand{\resp}{respectively\ }
\newcommand{\Subsection}{\subsection}
\providecommand{\nobreakdash}{\nobreak\hbox{-}}
\setlength{\emergencystretch}{2em}
\multlinegap0pt
\usepackage{amssymb}
\usepackage[shortlabels]{enumitem}
\usepackage{xcolor}
\newtheorem{theorem}{Theorem}
\newtheorem{lemma}{Lemma}
\def\bC{\mathbb{C}} 
\def\bR{\mathbb{R}}
\def\cU{\mathcal{U}}
\def\half{\tfrac{1}{2}}
\def\tPhi{\tilde\Phi}
\def\tRe{\mathrm{Re}} \def\tRic{\mathrm{Ric}}
\def\tand{\quad\hbox{and}\quad}
\title{Extension of Hodge norms at infinity}
\author{Colleen Robles}
\date{Source: arXiv:2302.04014v2 (CC BY 4.0)}
\begin{document}
\maketitle

\noindent\emph{Source and licence.} Extension of Hodge norms at infinity. Version arXiv:2302.04014v2 (CC BY 4.0), distributed by arXiv under the Creative Commons Attribution 4.0 International licence, http://creativecommons.org/licenses/by/4.0/, as recorded in the arXiv licence metadata for this exact version and shown on the abstract page https://arxiv.org/abs/2302.04014v2. Full licence notice: \texttt{licenses/arxiv-2302.04014-CC-BY-4.0.txt}.\par\smallskip

\noindent\textit{Editorial scope.} Counted unit: the article's theorem T:descent (source0.tex lines 876--878). The article's complete original proof of it is delivered with it (source0.tex lines 880--904, one \texttt{proof} environment of 25 lines). No mathematics is written, repaired or re-derived: the statement and the proof are the article's own lines, in the article's own notation, and no retained line is substituted or re-flowed.  Retained uncounted as the source-local context the proof needs: lines 869--873; lines 906--914.  Nothing was omitted from any retained span, so the package carries no omission record.  No retained line contains a citation, so the wrapper carries no bibliography and no bibliography entry.  No retained line contains a figure environment, an image-inclusion command or drawing code, so no image is delivered and the package declares no asset dependency.  Editorial formatting only, in addition: an AMS \texttt{amsart} preamble that restates the article's own declarations and macros used by the retained lines, namely the environments \texttt{theorem}, \texttt{lemma} on separate counters, together with the standard packages those lines need.  The retained environments print in this document as Theorem 1, Lemma 1 in order of appearance, whereas the article prints the same items with the numbering of its own full document.  The complete native source of the article as distributed in the arXiv e-print for this exact version is bundled unchanged as \texttt{sources/137109/source0.tex}.
\begin{equation}\label{E:tildeh}
  \tilde h \ = \ \half\,Q(\eta_0 \,,\,\overline{\lambda\,\eta_\infty})
  \,+\, \half\,Q(\overline{\eta_0} \,,\,\lambda\,\eta_\infty) 
  \ = \ \tRe\,Q(\eta_0 \,,\,\overline{\lambda\,\eta_\infty}) \,.
\end{equation}

\begin{lemma} \label{L:char}
There exists a character $\chi(\gamma) \in S^1 \subset \bC$ so that
\[
  \gamma \cdot \eta_0(u) \ = \ \chi(\gamma)\,\eta_0(\gamma \cdot u)
  \tand 
  \gamma \cdot \eta_\infty(u) \ = \ 
  \chi(\gamma)\,\eta_\infty(\gamma \cdot u) \,.
\] 
\end{lemma}

\begin{theorem} \label{T:descent}
The function $\tilde h : \cU \ \to \ \bR$ descends to $h:B \cap X \to \bR$.
\end{theorem}

\begin{proof}
To prove the theorem, we need to show that 
\begin{equation}\label{E:descent} 
  Q\left(\eta_0(u) , \overline{\eta_\infty(u)}\right) \ = \ Q\left(\eta_0(\gamma\cdot u) , \overline{\eta_\infty(\gamma \cdot u)}\right)
\end{equation}
for all $u \in \cU$ and $\gamma \in \pi_1(B \cap X)$, cf.~\eqref{E:tildeh}.  Note that $\pi_1(B \cap X)$ acts on both $\cU$ and $H_\bC$.  Given $u \in \cU$ and $\gamma \in \pi_1(B \cap X)$, we have
\[
  \gamma \cdot F^p(\tPhi(u)) \ = \ F^p(\tPhi(\gamma\cdot u)) \,.
\]
If $\eta_j(u) \in F^p(\tPhi(u))$, then it is likewise the case that both
\begin{equation}\label{E:gamma.eta}
  \gamma\cdot \eta_j(u)\,,\, \eta_j(\gamma \cdot u) 
  \ \in \ 
  F^p(\tPhi(\gamma\cdot u)) \,.
\end{equation}
However in general it need not be the case that $\gamma\cdot \eta_j(u) =  \eta_j(\gamma \cdot u)$, or even that they be linearly dependent.  Nonetheless, Lemma \ref{L:char} below does hold, and (keeping in mind that $\gamma \in \Gamma$ preserves the polarization $Q$) implies
\[
  Q\left( \eta_0(u) , \overline{\eta_\infty(u)} \right) 
  \ = \ 
  Q\left( \gamma \cdot \eta_0(u) , \gamma \cdot \overline{\eta_\infty(u)} \right) 
  \ = \ 
  Q\left( \eta_0(\gamma \cdot u) , \overline{\eta_\infty(\gamma \cdot u)} \right) \,,
\]
establishing the desired \eqref{E:descent}.
\end{proof}

\end{document}
